%-------------------------
% Resume in Latex
% Author : Sourabh Bajaj
% License : MIT
%------------------------

\documentclass[letterpaper,11pt]{article}

\usepackage{latexsym}
\usepackage[empty]{fullpage}
\usepackage{titlesec}
\usepackage{marvosym}
\usepackage[usenames,dvipsnames]{color}
\usepackage{enumitem}
\usepackage[pdftex]{hyperref}
\usepackage{fancyhdr}

\pagestyle{fancy}
\fancyhf{}
\fancyfoot{}
\renewcommand{\headrulewidth}{0pt}
\renewcommand{\footrulewidth}{0pt}

% Adjust margins
\addtolength{\oddsidemargin}{-0.375in}
\addtolength{\evensidemargin}{-0.375in}
\addtolength{\textwidth}{1in}
\addtolength{\topmargin}{-.5in}
\addtolength{\textheight}{1.0in}

\urlstyle{same}
\raggedbottom
\raggedright
\setlength{\tabcolsep}{0in}

% Sections formatting
\titleformat{\section}{
  \vspace{-4pt}\scshape\raggedright\large
}{}{0em}{}[\color{black}\titlerule \vspace{-5pt}]

%-------------------------
% Custom commands
\newcommand{\resumeItem}[2]{
  \item\small{
    \textbf{#1}{: #2 \vspace{-3pt}}
  }
}

\newcommand{\resumeSubheading}[4]{
  \vspace{-2pt}\item
    \begin{tabular*}{0.97\textwidth}{l@{\extracolsep{\fill}}r}
      \textbf{#1} & #2 \\
      \textit{\small#3} & \textit{\small #4} \\
    \end{tabular*}\vspace{-6pt}
}

\newcommand{\resumeSubItem}[2]{\resumeItem{#1}{#2}\vspace{-5pt}}

\renewcommand{\labelitemii}{$\circ$}

\newcommand{\resumeSubHeadingListStart}{\begin{itemize}[leftmargin=*]}
\newcommand{\resumeSubHeadingListEnd}{\end{itemize}\vspace{-8pt}}
\newcommand{\resumeItemListStart}{\begin{itemize}}
\newcommand{\resumeItemListEnd}{\end{itemize}\vspace{-6pt}}


%-------------------------------------------
%%%%%%  CV STARTS HERE  %%%%%%%%%%%%%%%%%%%%%%%%%%%%


\begin{document}

%----------HEADING-----------------
\begin{tabular*}{\textwidth}{l@{\extracolsep{\fill}}r}
  \textbf{\Large Sırrı Kaan Oğuzkan} & E-Mail : \href{mailto:kaan.oguzkan@ug.bilkent.edu.tr}{kaan.oguzkan@ug.bilkent.edu.tr}\\
  \href{https://www.linkedin.com/in/kaan-oguzkan/}{{linkedin.com/in/kaan-oguzkan}} & \href{https://kaanoguzkan.com}{{kaanoguzkan.com}} \\
  \href{https://github.com/kaanoguzkan}{{github.com/kaanoguzkan}} & \href{https://orcid.org/0009-0000-3272-7333}{{ORCID: 0009-0000-3272-7333}} \\
\end{tabular*}

\section{Ausbildung}
  \resumeSubHeadingListStart
    \resumeSubheading
      {Bilkent Universität}{Ankara, Türkei}
      {B.Sc. Informatik (Computer Engineering)}{Sep. 2023 -- heute}
  \resumeSubHeadingListEnd
%-----------EXPERIENCE-----------------
\section{Berufserfahrung}
  \resumeSubHeadingListStart

    \resumeSubheading
  {JotForm}{Ankara, Türkei}
  {Praktikant Software-Entwicklung}{Juni 2026 - Aug. 2026}
  \resumeItemListStart
      \resumeItem{Element-Registry-Architektur}
        {Plugin-Contract-Architektur für einen Drag-and-Drop-E-Mail-Builder (\textbf{React 18}, \textbf{TypeScript}, \textbf{Redux}) entworfen, bei der \textbf{20} Content-Block-Typen über Canvas, Panel, Vorschau und Exporter hinweg eine einzige typisierte Schnittstelle teilen; neue Blöcke erforderten so \textbf{keine} Änderung bestehender Konsumenten.}
      \resumeItem{Bulk-Send-Pipeline}
        {Asynchrone Bulk-E-Mail-Pipeline (PHP, MySQL, Redis) gebaut, mit einem Cron-Worker, der über gestückelte SMTP-Sitzungen (\textbf{25} Empfänger/Batch) versendet, und einem timeoutbasierten Reclaim-Mechanismus (\textbf{900 s}) für hängende Batches.}
    \resumeItemListEnd

    \resumeSubheading
  {Alkan Lab für Bioinformatik und Computational Genomics}{Ankara, Türkei}
  {Studentische Hilfskraft (ehrenamtlich)}{Feb. 2026 - heute}
  \resumeItemListStart
      \resumeItem{CuCLARK GPU- \& CUDA-Modernisierung}
        {CuCLARK (GPU-Metagenom-Klassifikator, urspr. CUDA 7.5 / Single-Arch) geforkt und modernisiert; auf \textbf{CUDA 11.8} aktualisiert, mit Multi-Arch-Support für \textbf{10} GPU-Generationen (Kepler \texttt{sm\_35} bis Hopper \texttt{sm\_90}, u.\,a. K80, V100, A100, RTX 30xx/40xx, H100).}
      \resumeItem{Build-System \& Containerisierung}
        {Legacy-\texttt{install.sh} durch einen \textbf{CMake-3.20}-Build (C++17, CUDA 14, optional OpenMP) ersetzt und ein Multi-Arch-\texttt{Dockerfile} verfasst (veröffentlicht als \texttt{alkanlab/cuclark}); einheitliche CLI mit \textbf{6} Unterbefehlen ergänzt, die die ursprünglichen Skripte kapselt.}
      \resumeItem{Batch-Klassifikations-Pipeline}
        {Byte-Offset-Batching-Bug diagnostiziert, der bei kleinen Eingaben doppelte Read-Zuordnungen verursachte; List-Mode-Pipeline mit Resume, Crash-Recovery-Blacklist und Fortschrittsüberwachung gebaut, die \textbf{3000+} FASTA-Dateien in einem einzigen DB-Ladevorgang klassifiziert.}
    \resumeItemListEnd

    \resumeSubheading
  {Look \& Cash}{Ankara, Türkei}
  {Werkstudent Software-Entwicklung}{Feb. 2025 - Okt. 2025}
  \resumeItemListStart
      \resumeItem{Ticket-basiertes Support-System}
        {Ticket-Workflow für Web/Mobile entwickelt (Priorität, Kommentare, Anhänge, Status). \textbf{13} REST-Endpunkte, \textbf{4} Ticket-Zustände und \textbf{5} Admin-Aktionen ausgeliefert; durchschnittliche Bearbeitungszeit der Admins durch Massenaktionen und bessere Filter um \textbf{25\,\%} reduziert.}
      \resumeItem{Echtzeit-Support-Chat (Mobile)}
        {Echtzeit-Ticket-Chat auf Mobile mit Backend-Persistenz und State-Sync implementiert; \textbf{300 ms} p95-Zustellung und \textbf{99,5\,\%} Erfolgsrate bei \textbf{150} parallelen Sitzungen erreicht.}
      \resumeItem{Reporting \& Diagrammerstellung}
        {KPI-Reporting-APIs auf MongoDB (Tickets, Kampagnen, Scans, Benachrichtigungen) mit Datums- und Kampagnenfiltern entwickelt. \textbf{14} Dashboard-Datensätze für \textbf{4} Diagrammtypen erzeugt; Reportzeit durch Indizierung von \textbf{6,2 s} auf \textbf{1,9 s} reduziert.}
    \resumeItemListEnd

  \resumeSubHeadingListEnd


%-----------PROJECTS-----------------
\section{Projekte}
  \resumeSubHeadingListStart

\resumeSubItem{osguide.org — Interaktive Lernplattform für Betriebssysteme}
{Interaktive Plattform (\textbf{React}, \textbf{TypeScript}, \textbf{Vite}) entwickelt, die sich am CS342-Lehrplan der Bilkent orientiert.}
\resumeItemListStart
  \resumeItem{Von Konzept zu Simulator}
    {Pro Thema einen dreistufigen Lernfluss entworfen — Konzepterklärung, Animation und interaktiver Simulator — der zentrale Betriebssystemthemen wie Scheduling, Speicherverwaltung und Synchronisation abdeckt.}
\resumeItemListEnd

\resumeSubItem{NutriApp — Plattform für Mahlzeitenempfehlungen}
{Backend-Dienste für eine Meal-Planning-App im Amazon University Engagement Program entwickelt.}
\resumeItemListStart
  \resumeItem{Serverless-REST-APIs}
    {\textbf{16} REST-Endpunkte auf AWS Lambda implementiert; \textbf{People's Choice Award} gewonnen, \textbf{Platz 1 von 7} Teams.}
\resumeItemListEnd

  \resumeSubHeadingListEnd


%-----------SKILLS-----------------
\section{Kenntnisse}
 \resumeSubHeadingListStart
    \item{
      \textbf{Programmiersprachen}{: C++, TypeScript, JavaScript, PHP, Java}
    }
    \item{
      \textbf{Frameworks \& Tools}{: React, Redux, Node.js, AWS Lambda, CUDA, Docker, Git}
    }
    \item{
      \textbf{Datenbanken}{: MongoDB, DynamoDB, MySQL, Redis}
    }

 \resumeSubHeadingListEnd


\end{document}
